\documentclass{article}

% ---------------- PACOTES ----------------
\usepackage[utf8]{inputenc}
\usepackage[T1]{fontenc}
\usepackage[brazilian]{babel}
\usepackage{array}
\usepackage{float}
\usepackage{listings}
\usepackage{amsmath}
\usepackage{geometry}

\geometry{a4paper, margin=2.5cm}

% ---------------- CONFIGURAÇÃO DO CÓDIGO C ----------------
\lstset{
    basicstyle=\ttfamily\small,
    breaklines=true,
    frame=single,
    language=C
}

\begin{document}

% =========================================================
% ===================== PÁGINA 1 ===========================
% =========================================================

\begin{titlepage}
\centering

{\Large UNIVERSIDADE PÚBLICA DE PERNAMBUCO\par}

\vspace{2cm}

{\large ALUNOS:\par}

\vspace{0.5cm}

Alline Heveline Bezerra\\
Carmen Lucia Layme\\
Hugo Smith Branquinho\\
Juan Kalebe Andrade\\
Kauwane Fidelis De Souza\\
Lucas Heitor Rodrigues\\
Rhelter Santana de Freitas\\
Sandro Goncalves da Silva Filho\\
Victor de Lima Melo\\
Yasmim Silva Maciel\\

\vfill

{\LARGE \textbf{FUNDAMENTOS DA LINGUAGEM C: ESTUDO SOBRE TIPOS DE DADOS}\par}

\vfill

\end{titlepage}

% =========================================================
% ===================== PÁGINA 2 ===========================
% =========================================================

\begin{titlepage}

\centering

{\Large UNIVERSIDADE PÚBLICA DE PERNAMBUCO\par}

\vspace{1.5cm}

{\large ALUNOS:\par}

\vspace{0.5cm}

Alline Heveline Bezerra\\
Carmen Lucia Layme\\
Hugo Smith Branquinho\\
Juan Kalebe Andrade\\
Kauwane Fidelis De Souza\\
Lucas Heitor Rodrigues\\
Rhelter Santana de Freitas\\
Sandro Goncalves da Silva Filho\\
Victor de Lima Melo\\
Yasmim Silva Maciel\\

\vfill

{\LARGE \textbf{FUNDAMENTOS DA LINGUAGEM C: ESTUDO SOBRE TIPOS DE DADOS}\par}

\vspace{1cm}

\begin{flushright}
Trabalho apresentado à disciplina de Programação como requisito parcial para obtenção de nota. Professor: Dr. Ruben Carlo Benante
\end{flushright}

\vfill

\end{titlepage}

% =========================================================
% ===================== CONTEÚDO ===========================
% =========================================================

\section{Fundamentos da Linguagem C}

A linguagem C é considerada uma linguagem de "médio nível" porque permite a manipulação direta da memória (ponteiros) ao mesmo tempo em que oferece estruturas de alto nível. Para dominar C, o primeiro passo é entender como os dados são armazenados e gerenciados.

\subsection{Tipos Básicos de Dados}

Em C, os tipos de dados definem o tamanho e o tipo de informação que uma variável pode armazenar. Os tipos fundamentais são:

\begin{table}[h]
\centering
\begin{tabular}{|l|l|l|l|}
\hline
\textbf{Tipo} & \textbf{Descrição} & \textbf{Tamanho} & \textbf{Intervalo Aproximado} \\ \hline
\texttt{char} & Caractere individual & 1 byte & -128 a 127 \\ \hline
\texttt{int} & Números inteiros & 4 bytes & $\pm$ 2,1 bilhões \\ \hline
\texttt{float} & Ponto flutuante & 4 bytes & 6 a 7 casas decimais \\ \hline
\texttt{double} & Precisão dupla & 8 bytes & 15 casas decimais \\ \hline
\texttt{void} & Tipo vazio & - & - \\ \hline
\end{tabular}
\caption{Tipos de dados fundamentais em C.}
\end{table}

\subsubsection*{Modificadores de Tipo}

\begin{itemize}
    \item \textbf{signed / unsigned:} Define se o número aceita sinais negativos ou apenas positivos.
    \item \textbf{short / long:} Altera o espaço de memória reservado para o dado.
\end{itemize}

\subsection{Declaração de Variáveis}

Uma variável é um espaço nomeado na memória. A declaração segue a sintaxe:
\texttt{tipo nome\_da\_variavel = valor\_inicial;}

\textbf{Regras de Ouro:}
\begin{itemize}
    \item \textbf{Tipagem Estática:} Você deve declarar o tipo da variável antes de usá-la.
    \item \textbf{Case Sensitive:} \texttt{idade} é diferente de \texttt{Idade}.
    \item \textbf{Nomes Válidos:} Devem começar com letra ou underscore (\_), nunca com números.
\end{itemize}

\begin{lstlisting}[caption={Exemplo de declaracao de variaveis}]
#include <stdio.h>

int main() {
    int idade = 25;           
    float altura = 1.75;      
    char inicial = 'G';       
    
    printf("Idade: %d, Altura: %.2f, Inicial: %c\n", idade, altura, inicial);
    return 0;
}
\end{lstlisting}

\subsection{Escopo de Variáveis}

O escopo define onde uma variável pode ser acessada dentro do programa.

\begin{itemize}
    \item \textbf{Variáveis Locais:} Declaradas dentro de uma função ou bloco \texttt{\{ \}}. Só existem durante a execução daquele bloco.
    \item \textbf{Variáveis Globais:} Declaradas fora de qualquer função. Podem ser acessadas e modificadas por todo o programa.
\end{itemize}

\begin{lstlisting}[caption={Exemplo de Escopo}]
#include <stdio.h>

int global_score = 100;

void test_function() {
    int local_val = 10;
    printf("Global: %d, Local: %d\n", global_score, local_val);
}

int main() {
    test_function();
    return 0;
}
\end{lstlisting}

\end{document}